\section{Phenomenological ansatz}
\label{sec:ansatz}

\subsection{Selected fitting functions}
\label{sec:hubble-ansatz}

Let $\tP = \sqrt{\hbar G/c^5} \approx 5.391\times 10^{-44}\,\mathrm{s}$ denote the Planck time, used only to make the argument of the logarithm dimensionless.
The function studied in this paper is $\xi(t)$ in \cref{eq:xi}.
For the Hubble comparison the working convention, taken from the figure and forecast notebooks rather than from the unscaled equation printed in the submitted text, is
\begin{equation}
  H(t) = H_\infty + \beta\,\frac{10^{5}}{\ln^2(t/\tP)}
       = H_\infty + \beta\,10^{5}\,\xi(t)\,.
  \label{eq:hubble}
\end{equation}
Here $H$ and $H_\infty$ are expressed in $\mathrm{km\,s^{-1}\,Mpc^{-1}}$.
The regressor $10^{5}\xi(t)$ is dimensionless and of order five over the redshift range discussed below, so $\beta$ carries the same units as $H$.
The submitted manuscript printed $H(t) = H_\infty + \beta/\ln^2(t/\tP)$ and quoted $\beta = -6.24$.
Substituting that pair into the unscaled expression changes $H$ by only $\sim 10^{-4}\,\mathrm{km\,s^{-1}\,Mpc^{-1}}$.
The scaled expression~\eqref{eq:hubble} is the one that the forecast notebook evaluates.
With the quoted pair $(H_\infty,\beta) = (104.5,-6.24)$ and the submitted leverage value $10^{5}\xi \approx 5.19$ at $z \approx 2$,
\begin{equation}
  H \approx 104.5 - 6.24\times 5.19 = 72.1\,,
\end{equation}
in $\mathrm{km\,s^{-1}\,Mpc^{-1}}$, which matches the value $H(z=2) = 72.1$ quoted in the submitted forecast.
Thus $H_\infty = 104.5\,\mathrm{km\,s^{-1}\,Mpc^{-1}}$ is the intercept at $\xi\to 0$, corresponding to $t\to\infty$ inside this fitting function.
The values near $71$--$72\,\mathrm{km\,s^{-1}\,Mpc^{-1}}$ are evaluations of the same curve at finite cosmic time.
They are not a second estimate of the intercept.
Because $\beta<0$, $H(t)$ in \cref{eq:hubble} increases toward $H_\infty$ as $t$ grows.
The submitted statement $H_\infty < H_0$ does not describe this fit.

A cosmological term of the same shape was written in the submitted manuscript as
\begin{equation}
  \Lambda(t) = \Lambda_\infty + \frac{\gamma}{\ln^2(t/\tP)}\,.
  \label{eq:lambda}
\end{equation}
In the Friedmann equation a cosmological term that dominates the expansion satisfies $H = \sqrt{\Lambda/3}$ only in de~Sitter spacetime, and only when $\Lambda$ is expressed in units of time$^{-2}$ (with $c=1$ in the usual relativistic convention).
The intercept $H_\infty$ in \cref{eq:hubble} is a fit coefficient in $\mathrm{km\,s^{-1}\,Mpc^{-1}}$.
It is not converted to a $\Lambda$ in this revision, and the identification $H_\infty = \sqrt{\Lambda_\infty/3}$ is not used.
Whether $\gamma$ in \cref{eq:lambda} includes the same $10^{5}$ normalization as $\beta$ is not fixed by the Hubble notebook.
\revtodo{Adopt one unit convention for $(\Lambda_\infty,\gamma)$ and state whether $\gamma$ multiplies $\xi$ or $10^{5}\xi$. Do not introduce a new fitted value until that convention is chosen.}

\subsection{Fine-structure parametrization}
\label{sec:alpha-ansatz}

The quasar comparison uses a second parametrization, already written in the submitted manuscript,
\begin{equation}
  \frac{\Delta\alpha}{\alpha}(t)
  = \Gamma\times 10^{-5}\,\big[\xi(t)-\xi(t_{\mathrm{now}})\big]\,,
  \label{eq:alpha}
\end{equation}
with $\xi$ the unscaled function~\eqref{eq:xi} and with $\Gamma$ quoted in units of $10^{-5}$.
Differentiating with respect to cosmic time gives
\begin{equation}
  \left.\frac{\dot\alpha}{\alpha}\right|_{t}
  = \Gamma\times 10^{-5}\times\frac{-2}{t\,\ln^3(t/\tP)}\,.
  \label{eq:alphadot}
\end{equation}
The submitted manuscript evaluates the prefactor at the present epoch as $5.25\times 10^{-22}\,\mathrm{yr}^{-1}$ per unit $\Gamma$.
That arithmetic follows from $\ln(t_{\mathrm{now}}/\tP)\approx 140$ for a present age of order $10^{10}\,\mathrm{yr}$; it is a property of \cref{eq:alphadot}, and it does not depend on a quasar fit.
The Rosenband et al.\ bound $|\dot\alpha/\alpha| < 1\times 10^{-17}\,\mathrm{yr}^{-1}$~\citep{Rosenband2008} then translates, inside this ansatz, to $|\Gamma|\lesssim 1.9\times 10^{4}$.
This is a conditional laboratory translation.
It is not evidence that a nonzero $\Gamma$ has been measured, and the Hubble coefficient $\beta$ is not the same parameter as $\Gamma$.

\subsection{What the ansatz does not assert}
\label{sec:not-derived}

The shape $1/\ln^2(t/\tP)$ was selected because it varies slowly for $t\gg\tP$.
A power-law drift that is large at $z\sim 2$ is difficult to reconcile with a tight present-day laboratory bound, which is a reason to consider slow functions.
That observation does not select the exponent $2$, and it does not rule out other slow functions such as $1/\ln^p(t/\tP)$ for $p\neq 2$.
No spectral density, symplectic map, or lattice Hamiltonian is used to derive \cref{eq:hubble,eq:alpha}.
The cosmic time $t$ that enters $\xi(t)$ must be supplied by an external $t(z)$.
The forecast notebook uses a flat $\Lambda$CDM age integral for that conversion.
Using a $\Lambda$CDM clock to test a proposed departure from $\Lambda$CDM is an assumption of the comparison and is stated as such in \cref{sec:methods}.
